\scanpage{86}
\mspage{3}
% archive sheet number: 85
1.11. Le groupe ${\unsure{\Phi|\varphi}}^i_\nu$ peut être appelé \uline{groupe des cycles évanescents en dimension $i$} (pour $X/S$\inserted{et les coeff.\ $\Lambda_\nu$}) et la suite exacte (1.9) la \uline{suite exacte des cycles évanescents}. Les cycles évanescents apparaissent ici comme mesurant ``l'écart'' entre la cohomologie de $X_s$ et de $X_{\bar{\unsure{\eta|y}}}$, fibres géométriques spéciale et générale.

Passant à la limite \inserted{projective} sur $\nu$, on trouve (les groupes de la suite (1.9) étant \uline{finis} \unsure{$\ell$)|(!)} par le th. de finitude)
\[
\cdots\longrightarrow\Phi^i\longrightarrow H^i(X_s,\mathbb Z_\ell)\longrightarrow H^i(X_{\bar\eta},\mathbb Z_\ell)\longrightarrow\Phi^{i+1}\longrightarrow\cdots,
\tag*{(1.\deleted{\unsure{9}}\deleted{\unsure{9}}9 bis)}
\]
avec
\[
\Phi^i=\varprojlim_\nu\Phi^i_\nu=\varprojlim_\nu\varinjlim_\alpha H^i_{X_{\alpha0}}(X_\alpha,\Lambda_\nu).
\tag*{(1.10 bis)}
\]

\deleted{1.12.} On appelle aussi (1.9 bis) suite exacte des cycles évanescents pour les coeff.\ $\mathbb Z_\ell$, et $\Phi^i$ le \unsure{groupe} des cycles évanescents en dim.\ $i$ pour les coeff.\ $\mathbb Z_\ell$.

1.12. \uline{Passage du local au global} Pour calculer les $\Phi^i_\nu$, \deleted{il y a lieu} utilisant la \deleted{\unsure{repr}}\unsure{expression} (1.10), on regarde pour chaque $\alpha$ la suite spectrale
\[
H^*_{X_{\alpha0}}(X_\alpha,\Lambda_\nu)\Longleftarrow E_2^{pq}=H^p(X_{\alpha0},\underline H^q_{X_{\alpha0}}(\Lambda_{\nu X_\alpha}))
\tag*{(1.13)}
\]
Ces suites spectrales, pour $\alpha$ variable, forment un système inductif, et donnent à la limite, compte tenu que $X_{\alpha0}\longrightarrow X_0$ induit une équivalence de \unsure{sites étales}
